\documentclass[conference]{IEEEtran}
\usepackage[utf8]{inputenc}
\usepackage[T1]{fontenc}
\usepackage[spanish,es-tabla]{babel}
\usepackage{graphicx}
\usepackage{array}
\usepackage{booktabs}
\usepackage{amsmath}
\usepackage{url}
\usepackage{cite}

\begin{document}

\title{\begin{tabular}{@{}m{2.3cm}@{\hspace{0.3cm}}m{\dimexpr\textwidth-3.0cm\relax}@{}}

\includegraphics[width=2.1cm]{figuras/logo_uni.png} &
\centering\arraybackslash Extracci\'on autom\'atica de resultados electorales
a partir de actas de escrutinio de las Elecciones Generales del Per\'u 2026
\end{tabular}}

\author{
\IEEEauthorblockN{Josemanuel Rossy Ca\~nari Palante, Kenny Asto Hinostroza,
Melissa Dessire Aylas Barranca, Carlos P\'erez P\'erez}
\IEEEauthorblockA{Maestr\'ia en Inteligencia Artificial --- Visi\'on por Computador, Grupo 3\\
Universidad Nacional de Ingenier\'ia, Lima, Per\'u}
}

\maketitle

\begin{abstract}
En el Per\'u, los resultados electorales oficiales se publican mediante actas
de escrutinio, muchas de ellas escaneadas. Esta situaci\'on dificulta el
procesamiento autom\'atico de la informaci\'on a gran escala. Se presenta un
sistema completo que descarga las actas del portal de la Oficina Nacional de
Procesos Electorales (ONPE), ubica 46 campos de cada documento mediante una
plantilla y marcas de referencia impresas, y reconoce los conteos manuscritos
con reconocimiento \'optico de caracteres (OCR) limitado a d\'igitos. La
evaluaci\'on no requiri\'o etiquetado manual: se usaron como referencia los
resultados digitados y publicados por la ONPE. El conjunto de prueba estuvo
formado por una muestra aleatoria de 100 actas de 23 departamentos. El 24\% de las actas fueron documentos digitales con
capa de texto extra\'ible y alcanzaron 95.1\% de exactitud por campo. En las 76
actas manuscritas escaneadas (3,268 campos), el sistema logr\'o 59.6\% de
exactitud por campo e \'indice de error por car\'acter (CER) de 0.401. En las
mismas actas, un Modelo de Lenguaje Grande (LLM) multimodal obtuvo 48.4\%. Los resultados
muestran que un sistema especializado y dividido por etapas es una base \'util
para mejorar el reconocimiento de d\'igitos manuscritos en documentos
electorales.
\end{abstract}

\begin{IEEEkeywords}
an\'alisis de documentos, reconocimiento \'optico de caracteres,
reconocimiento de d\'igitos manuscritos, documentos electorales, alineaci\'on
de plantillas
\end{IEEEkeywords}

\section{Introducci\'on}
Las actas de escrutinio contienen los resultados oficiales de cada mesa de
sufragio en el Per\'u. Aunque la ONPE publica estos documentos en formato
digital, gran parte de la informaci\'on est\'a contenida en formularios
escaneados donde los miembros de mesa escribieron los votos a mano. Revisar
estas actas de manera manual es lento, puede generar errores y no es viable
para las aproximadamente 88,000 mesas de una elecci\'on general. Esto limita
la auditor\'ia autom\'atica y la verificaci\'on independiente de los resultados.

Este trabajo responde a la siguiente pregunta: \emph{\textquestiondown qu\'e
tan preciso puede ser un sistema tradicional de visi\'on por computador al
extraer resultados de actas peruanas seleccionadas a escala nacional y qu\'e
implica ello para un sistema de uso real?} Para ello, se construy\'o y evalu\'o
un proceso de cinco etapas: descarga, preparaci\'on de la imagen, ubicaci\'on de
campos, reconocimiento de d\'igitos y generaci\'on de una salida estructurada.
El estudio se realiz\'o con actas de la primera vuelta de las Elecciones
Generales 2026. Las principales contribuciones son las siguientes:

\begin{enumerate}
\item \textbf{Una evaluaci\'on reproducible sin etiquetado manual.} El portal
de la ONPE publica tanto la imagen del acta como los conteos digitados
oficialmente. Estos valores se utilizaron como referencia para evaluar el
sistema. As\'i, cada acta aporta alrededor de 43 pares de campo y valor ya
verificados. La muestra de prueba se seleccion\'o aleatoriamente a nivel
nacional con una semilla publicada.

\item \textbf{Una plantilla alineada con marcas impresas del formulario.} El
acta tiene un dise\~no fijo y contiene cuadrados negros en sus bordes. Estas
marcas permiten corregir desplazamientos y peque\~nas deformaciones producidas
durante el escaneo. Esta correcci\'on mejor\'o la exactitud en 8.7 puntos
porcentuales frente al uso de coordenadas fijas sin alineaci\'on.

\item \textbf{La identificaci\'on de dos tipos de actas.} El 24\% de la muestra
corresponde a actas generadas digitalmente, cuya capa de texto puede extraerse
directamente y alcanza 95.1\% de exactitud por campo. El 76\% restante son
actas manuscritas escaneadas, que requieren procesamiento de imagen. Por ello,
un sistema de uso real debe identificar primero el tipo de documento y aplicar
el proceso adecuado.

\item \textbf{Una medida clara del rendimiento del OCR general.} En las actas
manuscritas, el OCR limitado a d\'igitos alcanz\'o 59.6\% de exactitud por
campo. Este resultado muestra el l\'imite de usar un reconocedor general sin
entrenarlo con este tipo de documento y sustenta la necesidad de entrenar un
modelo especializado.

\item \textbf{Una comparaci\'on con un modelo de lenguaje multimodal.} Con las
mismas actas y condiciones de evaluaci\'on, un modelo multimodal que recibe la
p\'agina completa y dos ejemplos resueltos obtuvo 48.4\%, frente a 60.9\% del
sistema por etapas. El modelo multimodal facilita una primera aproximaci\'on,
pero no supera al m\'etodo especializado en la lectura de d\'igitos manuscritos.
\end{enumerate}

El c\'odigo, los resultados por campo y la semilla usada para el muestreo son
p\'ublicos.\footnote{\url{https://github.com/carlosperez100/onpe_actas}}

\section{Trabajos relacionados}
\textbf{Detecci\'on de objetos.} RT-DETR \cite{rtdetr} es un modelo de detecci\'on
de objetos en tiempo real que, seg\'un sus autores, supera a modelos YOLO de
tama\~no similar en exactitud y velocidad. En actas de escrutinio, donde las
casillas est\'an juntas y pueden estar cubiertas por sellos o firmas, los
detectores aprendidos podr\'ian ser una alternativa para localizar los campos.
En este trabajo, esa localizaci\'on se resuelve mediante una plantilla alineada
con el documento. En el futuro, esta plantilla tambi\'en puede generar
autom\'aticamente etiquetas para entrenar detectores como YOLOv11 o RT-DETR.

\textbf{Reconocimiento de texto.} TrOCR \cite{trocr} emplea redes neuronales
de tipo \emph{transformer} para el reconocimiento \'optico de caracteres.
DTrOCR \cite{dtrocr} propone una variante orientada al reconocimiento de texto
impreso, de escenas y manuscrito. Estos trabajos respaldan el uso de modelos
especializados para mejorar la lectura de d\'igitos manuscritos. En nuestro
caso, es importante evitar que un modelo modifique un n\'umero para convertirlo
en una secuencia aparentemente m\'as probable, ya que cada valor representa un
conteo electoral.

\textbf{An\'alisis integral de documentos.} OmniParser \cite{omniparser}
integra la ubicaci\'on de texto, la extracci\'on de datos importantes y el
reconocimiento de tablas. Donut \cite{donut}, en cambio, intenta generar la
estructura del documento directamente desde la imagen. Los resultados de estos
trabajos muestran que los documentos con tablas extensas pueden ser dif\'iciles
de procesar en un solo paso. Por ello, se mantiene un proceso por etapas, ya
que permite revisar y medir con claridad el desempe\~no de cada parte del
sistema. Los m\'etodos integrales se consideran una alternativa para trabajos
posteriores.

Los trabajos revisados no eval\'uan espec\'ificamente d\'igitos manuscritos en
actas electorales peruanas con sellos y firmas. Esta es la necesidad pr\'actica
que aborda el presente estudio.

\section{Conjunto de datos}
\subsection{Fuente y descarga}
Todos los documentos proceden del portal oficial de resultados de la ONPE para
la primera vuelta de las Elecciones Generales 2026, realizada el 12 de abril de
2026. El portal permite consultar las actas de una mesa, su estado de
procesamiento y los archivos asociados. El acta de escrutinio se descarga en
formato PDF desde una direcci\'on temporal proporcionada por el portal. Las
descargas se realizaron a una velocidad moderada y se registraron en un archivo
que permite retomarlas en caso de interrupci\'on (\texttt{src/scraper/}).

\subsection{Dise\~no de la muestra}
El universo de estudio contiene aproximadamente 88,064 c\'odigos de mesa,
estimados mediante consultas a los servicios del portal. Se extrajo una muestra
aleatoria simple, sin reemplazo, de 100 actas con una semilla publicada (2026). La
muestra cubre 23 departamentos, identificados mediante los prefijos de
\emph{ubigeo}, que es el c\'odigo geogr\'afico est\'andar utilizado en el Per\'u.
La participaci\'on de Lima fue de 36\%, coherente con su mayor peso en el
electorado. Este dise\~no evita el sesgo de seleccionar mesas
con c\'odigos consecutivos de un solo distrito.

\subsection{Tipos de actas}
La revisi\'on del n\'umero de p\'aginas y de la capa de texto de los PDF permiti\'o
identificar dos tipos de documentos: \textbf{76 actas manuscritas}, que son
escaneos de una p\'agina sin texto extra\'ible, y \textbf{24 actas digitales
STAE}, que tienen dos p\'aginas, firma digital y conteos escritos con tipograf\'ia
digital. STAE significa \emph{Soluci\'on Tecnol\'ogica de Apoyo al Escrutinio},
una herramienta de la ONPE que ayuda a los miembros de mesa a registrar e
imprimir las actas. En este estudio, las actas STAE son los documentos
generados digitalmente con dicha herramienta.

La proporci\'on de actas digitales en la muestra es 24\%, con un intervalo de
confianza de Wilson del 95\% entre 16.7\% y 33.2\%. Este resultado confirma que
la poblaci\'on incluye ambos tipos de documentos, aunque la muestra no pretende
fijar su proporci\'on exacta. Los experimentos de visi\'on se realizaron con las
76 actas manuscritas; las actas digitales se procesaron mediante la extracci\'on
directa de la capa de texto (Secc.~\ref{sec:stae}).

\subsection{Valores de referencia}
Para cada mesa de la muestra se guardaron los conteos digitados oficialmente
por la ONPE: votos por organizaci\'on pol\'itica, votos en blanco, nulos e
impugnados, totales y n\'umero de electores. Las posiciones 04 y 14 del
formulario corresponden a partidos retirados; por ello est\'an vac\'ias en el
acta y no tienen valor en los datos oficiales. Se excluyeron de la evaluaci\'on
para evitar contar como aciertos coincidencias entre campos vac\'ios. En total,
cada acta aporta 43 campos evaluables, es decir, 3,268 campos en las actas
manuscritas.

\section{M\'etodo}
\subsection{Descripci\'on general}
El sistema consta de cinco etapas: (1) descarga del PDF y de los valores
oficiales de referencia; (2) preparaci\'on de la imagen; (3) ubicaci\'on de las
regiones de inter\'es; (4) reconocimiento de los valores; y (5) generaci\'on de
un archivo estructurado por acta. La preparaci\'on de la imagen incluye
convertir el PDF a imagen con una resoluci\'on de 300 p\'ixeles por pulgada,
corregir la inclinaci\'on, mejorar el contraste y reducir el ruido de escaneo.

\subsection{Ubicaci\'on de campos mediante plantilla y marcas de referencia}
\label{sec:registro}
El acta tiene un formato fijo. Por ello, se definieron una sola vez 46 regiones
en una acta de referencia: 38 celdas de votos, cuatro filas de resumen y cuatro
campos de encabezado. Las coordenadas se guardaron como proporciones del ancho
y alto de la p\'agina (Fig.~\ref{fig:overlay}).

El uso directo de coordenadas fijas falla cuando la imagen presenta m\'argenes
negros, desplazamientos o ligeras deformaciones producidas por el esc\'aner. En
algunas actas, el contenido se desplaz\'o, haciendo que una celda se recortara sobre la fila equivocada. Para
resolverlo, se usan los cuadrados negros impresos en los bordes del formulario
como marcas de referencia. Primero se detectan los cuadrados oscuros cercanos a
los bordes. Luego se comparan con 15 posiciones conocidas del acta de referencia
para calcular el desplazamiento m\'as consistente, ignorando marcas ausentes o
falsas detecciones. Finalmente, se aplica una correcci\'on geom\'etrica simple
a la p\'agina antes de recortar los campos.

\begin{figure}[t]
\centering
\includegraphics[width=\columnwidth]{figuras/plantilla_overlay.png}
\caption{Acta de referencia con las 46 regiones de la plantilla: 38 celdas de
votos (azul), campos de resumen y encabezado (rojo). Las coordenadas se aplican
despu\'es de alinear el documento con las marcas de referencia.}
\label{fig:overlay}
\end{figure}

\subsection{Reconocimiento de conteos manuscritos}
Cada celda se procesa con EasyOCR, configurado para reconocer solamente
d\'igitos. Antes del reconocimiento se aplican mejoras obtenidas a partir de
un an\'alisis de errores con 10 actas de prueba (Fig.~\ref{fig:limpieza}).
Primero, se elimina un margen interno para evitar que los bordes impresos de
la celda sean interpretados como d\'igitos. Despu\'es, se convierte la imagen a
blanco y negro con un umbral calculado para cada celda, de forma que se resalte
la tinta manuscrita. Tambi\'en se eliminan los separadores punteados impresos
en posiciones conocidas de la celda y las peque\~nas manchas producidas por el
escaneo. En los campos que contienen una casilla por d\'igito, cada casilla se
lee por separado. Si este proceso elimina trazos demasiado tenues, el sistema
vuelve a intentar la lectura con el recorte original.

\begin{figure}[t]
\centering
\includegraphics[width=\columnwidth]{figuras/limpieza_celdas.png}
\caption{Limpieza de celdas problem\'aticas (arriba: recorte original; abajo:
recorte procesado). Se eliminan bordes y separadores punteados que pod\'ian ser
confundidos con d\'igitos.}
\label{fig:limpieza}
\end{figure}

\subsection{Regla para las celdas vac\'ias}
Con frecuencia, los miembros de mesa dejan en blanco la celda de un partido
que no recibi\'o votos. En este formulario, una celda vac\'ia representa cero y
no necesariamente un error de OCR. Por ello, cuando no se reconoce ning\'un
valor en una celda de partido o en votos impugnados, el sistema asigna el valor
cero. Este aporte de conocimiento sobre el formulario se informa por separado,
para diferenciarlo del rendimiento propio del reconocimiento de imagen.

\subsection{Procesamiento de actas digitales}
\label{sec:stae}
En las actas digitales STAE, los resultados se extraen directamente de la capa
de texto del PDF. Cada nombre de partido aparece junto a su conteo y las
etiquetas de resumen aparecen antes de sus valores. Los nombres extra\'idos se
comparan con los nombres oficiales despu\'es de normalizar may\'usculas,
espacios y tildes. De esta forma, no es necesario procesar la imagen.

\section{Experimentos}
\subsection{M\'etricas y procedimiento estad\'istico}
La \textbf{exactitud por campo} es el porcentaje de campos cuyo valor le\'ido
coincide exactamente con el valor oficial. El \textbf{CER} (del ingl\'es
\emph{Character Error Rate}) es la tasa de error por car\'acter: compara las
cadenas de d\'igitos reconocidas con las cadenas correctas. Un valor menor de
CER indica menos errores.

Los campos de una misma acta pueden estar relacionados porque comparten la
misma calidad de escaneo y la misma escritura. Por esta raz\'on, los intervalos
de confianza se calcularon mediante remuestreo de actas completas, con 10,000
repeticiones. Los resultados por campo y las m\'etricas por acta est\'an
disponibles en el repositorio del proyecto.

\subsection{Evaluaci\'on gradual de mejoras en el piloto (10 actas)}
En un análisis inicial, se utiliz\'o mesas consecutivas de un solo distrito. Por tanto, no sirve para representar al pa\'is, pero s\'i permite medir el aporte de cada mejora incorporada al reconocimiento (Tabla~\ref{tab:pilot}).

\begin{table}[t]
\caption{Efecto acumulado de las mejoras en el piloto de 10 actas.}
\label{tab:pilot}
\centering
\begin{tabular}{lrr}
\toprule
Configuraci\'on & Exactitud (\%) & CER \\
\midrule
v1: OCR de d\'igitos sin limpieza & 43.95 & 0.524 \\
v2: + limpieza de tinta, bordes y casillas & 52.09 & 0.419 \\
v3: + eliminaci\'on de separadores y manchas & 52.79 & 0.476 \\
v4: + regla de celda vac\'ia = cero & 58.37 & 0.420 \\
\bottomrule
\end{tabular}
\end{table}

\subsection{Evaluaci\'on nacional (76 actas manuscritas)}
La Tabla~\ref{tab:national} presenta los resultados del sistema en la muestra
aleatoria nacional. La alineaci\'on mediante marcas de referencia result\'o
fundamental. Sin ella, las variaciones del encuadre afectaron gravemente a
algunas actas: la peor obtuvo 14\% de exactitud y un CER superior a 1. Con la
alineaci\'on, la exactitud aument\'o 8.7 puntos porcentuales y la peor acta
subi\'o de 14\% a 30\%. La regla de celdas vac\'ias equivalentes a cero a\~nadi\'o
21.6 puntos porcentuales, pues estos casos fueron frecuentes en la muestra
nacional.

\begin{table}[t]
\caption{Resultados en la muestra aleatoria nacional de 76 actas manuscritas.}
\label{tab:national}
\centering
\begin{tabular}{lrr}
\toprule
Configuraci\'on & Exactitud (\%) & CER \\
\midrule
Plantilla fija, OCR sin limpieza & 26.81 & 0.724 \\
Plantilla fija + regla de cero & 50.86 & 0.482 \\
Plantilla alineada, OCR sin limpieza & 38.00 & 0.617 \\
\textbf{Plantilla alineada + regla de cero} & \textbf{59.58} & \textbf{0.401} \\
\bottomrule
\end{tabular}
\end{table}

\textbf{Resultado principal.} El sistema obtuvo \textbf{59.58\% de exactitud
por campo}, con un intervalo de confianza del 95\% entre 56.61\% y 62.48\%.
El CER fue \textbf{0.401}, con intervalo de confianza del 95\% entre 0.360 y
0.444. La exactitud media por acta fue 59.6\%, con una desviaci\'on est\'andar
de 13.0 puntos y un rango entre 30\% y 84\%. Las diferencias restantes se
deben principalmente a la caligraf\'ia y a la calidad del escaneo, m\'as que a
la ubicaci\'on de los campos.

\textbf{Actas digitales.} La extracci\'on de la capa de texto en las 24 actas
STAE logr\'o \textbf{95.12\%} de exactitud (955 de 1,004 campos). Los errores
restantes se relacionan con la asociaci\'on entre nombres de partidos, no con
la lectura de los n\'umeros. La diferencia de 35.5 puntos porcentuales respecto
a las actas manuscritas confirma la importancia de identificar el tipo de
documento antes de procesarlo.

\subsection{Comparaci\'on con un modelo de lenguaje multimodal}
\label{sec:llm}
Como alternativa al proceso por etapas, se evalu\'o un modelo de lenguaje de
gran tama\~no con capacidad para analizar imagen y texto (LLM, por sus siglas en
ingl\'es). El modelo recibi\'o la imagen
completa del acta y dos ejemplos de actas con su resultado estructurado. Deb\'ia
devolver los conteos sin utilizar plantilla, alineaci\'on ni OCR
independiente. Se compararon ambos m\'etodos en las mismas 15 actas manuscritas. El modelo utilizado fue Gemini 2.5 Flash, en su capa gratuita,
consultado en julio de 2026. El c\'odigo est\'a disponible en
\texttt{src/recognition/llm\_ocr.py}.

\begin{table}[t]
\caption{Comparaci\'on entre el sistema por etapas y un modelo multimodal con
dos ejemplos, sobre las mismas 15 actas manuscritas.}
\label{tab:llm}
\centering
\begin{tabular}{lrr}
\toprule
M\'etodo & Exactitud (\%) & CER \\
\midrule
\textbf{Sistema por etapas (plantilla alineada + OCR)} & \textbf{60.93} & \textbf{0.395} \\
LLM multimodal con dos ejemplos & 48.37 & 0.591 \\
\bottomrule
\end{tabular}
\end{table}

El sistema por etapas super\'o al modelo multimodal en 12.6 puntos porcentuales.
El modelo identific\'o correctamente la estructura general del formulario, pero
tuvo m\'as dificultades para leer d\'igitos manuscritos en escaneos de baja
calidad. En las dos actas con peor resultado, obtuvo 20.9\% y 23.3\%, mientras
que el sistema por etapas obtuvo 34.9\% y 53.5\%, respectivamente. Esto se debe
a que el OCR recibe una celda ya aislada, mientras que el modelo multimodal
debe interpretar toda la p\'agina y su ruido visual. El modelo multimodal puede
ser \'util para preparar una primera versi\'on de etiquetas en un formato nuevo,
pero no sustituye a un reconocedor especializado en esta tarea.

\subsection{An\'alisis de errores}
Los errores se concentraron en cuatro situaciones: d\'igitos delgados que el
OCR no detecta; trazos manuscritos que invaden las zonas donde se eliminan los
separadores; confusiones entre formas similares, como 8 y 18 o 7 y 9; y campos
formados por varias casillas. Estos son principalmente errores de
reconocimiento y no de ubicaci\'on de los campos. Por ello, el siguiente paso
es entrenar un reconocedor espec\'ifico para estos d\'igitos y este formulario.

\section{Conclusiones y trabajo futuro}
El sistema propuesto alcanz\'o 59.6\% de exactitud por campo en actas
manuscritas seleccionadas a nivel nacional y 95.1\% en actas digitales. La
evaluaci\'on fue reproducible y no requiri\'o etiquetado manual, pues utiliz\'o
los resultados digitados oficialmente por la ONPE como referencia.

El estudio deja cuatro conclusiones principales. Primero, antes de recortar
campos con una plantilla es necesario alinear el documento con sus marcas
impresas de referencia. Segundo, la regla de que una celda vac\'ia equivale a
cero debe diferenciarse del rendimiento estrictamente visual del OCR. Tercero,
conviene identificar si el documento es manuscrito o digital antes de elegir
el m\'etodo de procesamiento. Cuarto, un modelo multimodal general requiere
menos preparaci\'on inicial, pero obtuvo 12.6 puntos porcentuales menos de
exactitud que el sistema especializado.

Como trabajo futuro, se propone entrenar un reconocedor basado en
\emph{transformers} con los pares de campo y valor obtenidos de las actas
(aproximadamente 43 por acta y 43,000 con 1,000 actas). Tambi\'en se propone
entrenar detectores de campos, como YOLOv11 y RT-DETR \cite{rtdetr}, usando
etiquetas creadas autom\'aticamente con la plantilla alineada; comparar el
sistema por etapas con m\'etodos integrales de an\'alisis de documentos
\cite{omniparser,donut}; y emplear el modelo multimodal como apoyo para generar
etiquetas iniciales que luego puedan ser revisadas por una persona.

\begin{thebibliography}{9}
\bibitem{rtdetr} Y.~Zhao, W.~Lv, S.~Xu, J.~Wei, G.~Wang, Q.~Dang, Y.~Liu,
and J.~Chen, ``DETRs beat YOLOs on real-time object detection,'' in
\emph{Proc. IEEE/CVF CVPR}, 2024, pp. 16965--16974.
\bibitem{dtrocr} M.~Fujitake, ``DTrOCR: Decoder-only transformer for optical
character recognition,'' in \emph{Proc. IEEE/CVF WACV}, 2024, pp. 8025--8035.
\bibitem{omniparser} J.~Wan, S.~Song, W.~Yu, Y.~Liu, W.~Cheng, F.~Huang,
X.~Bai, C.~Yao, and Z.~Yang, ``OmniParser: A unified framework for text
spotting, key information extraction and table recognition,'' in \emph{Proc.
IEEE/CVF CVPR}, 2024, pp. 15641--15653.
\bibitem{trocr} M.~Li, T.~Lv, J.~Chen, L.~Cui, Y.~Lu, D.~Florencio, C.~Zhang,
Z.~Li, and F.~Wei, ``TrOCR: Transformer-based optical character recognition
with pre-trained models,'' in \emph{Proc. AAAI}, 2023, pp. 13094--13102.
\bibitem{donut} G.~Kim, T.~Hong, M.~Yim, J.~Nam, J.~Park, J.~Yim, W.~Hwang,
S.~Yun, D.~Han, and S.~Park, ``OCR-free document understanding transformer,''
in \emph{Proc. ECCV}, 2022, pp. 498--517.
\end{thebibliography}

\end{document}
